\section{The finite character comparison}
\label{fcmp-section}

We now pass from character Selmer groups over $\mathbb Q$ to
the required eigenspace over an abelian CM field. The comparison
is exact at primes not dividing the field degree. At the other
odd primes an integral Fourier transform leaves a bounded
torsion error. In the actual family of fields those primes and
the necessary torsion exponents are fixed before the conductor
grows. A class-field calculation bounds their total contribution
by a fixed multiple of $N^2t\log N$.

\subsection{The local conditions and the parity}

Let $L/\mathbb Q$ be a finite abelian CM extension, let
$\Delta=\operatorname{Gal}(L/\mathbb Q)$, let
$d=[L:\mathbb Q]$, and let $j\in\Delta$ be complex
conjugation. Fix $N\ge3$ such that
$L\subset\mathbb Q(\mu_N)$, so $d\le\varphi(N)\le N$.
For an odd prime $p$ and $e\ge1$, put
$A_e=\mathbb Q_p/\mathbb Z_p(-e)$ and retain
$\mathcal S_p(L,e)$ from (MS1). Its $\Delta$-action
conjugates Galois arguments and acts on the coefficients.
Set
\[
 a_p(L,e)=\operatorname{length}_{\mathbb Z_p}
                  \mathcal S_p(L,e)^{\,j=(-1)^e}.
                                                               \tag{FC1}
\]
Thus (FC1) is exactly the complementary-parity summand.

Choose a finite extension $K/\mathbb Q_p$ containing the
values of every character of $\Delta$, and let $\mathcal O$
be its ring of integers, with ramification index $b$.
For $\theta\in\widehat\Delta$ define
\[
 \begin{split}
 A_{e,\theta}&=(K/\mathcal O)
                         (\theta^{-1}\varepsilon_p^{-e}),\\
 \mathcal V_p(e,\theta)
   &=\ker\left(H^1(\mathbb Q,A_{e,\theta})
              \longrightarrow\prod_{\ell\ {\rm prime}}
                         H^1(\mathbb Q_\ell,A_{e,\theta})\right),\\
 \ell_p(e,\theta)&=b^{-1}
       \operatorname{length}_{\mathcal O}\mathcal V_p(e,\theta).
 \end{split}                                                   \tag{FC2}
\]
We allow infinite length until finiteness is proved.

\begin{lemma}\label{fcmp-local}
For $e\ge1$, the strict Selmer group $\mathcal S_p(L,e)$
is the kernel of restriction to every finite decomposition
group. For a character $\theta$ of conductor dividing $N$,
the group $\mathcal V_p(e,\theta)$ is equivalently defined
by zero restriction to decomposition groups at primes
dividing $pN$, and zero restriction to inertia at all
other finite primes.
\end{lemma}
\begin{proof}
Let $w\nmid p$ be a place of $L$. Inertia acts trivially on
$A_e$, and arithmetic Frobenius acts by $q_w^{-e}$, where
$q_w>1$ is the residue-field cardinality. A class trivial
on inertia comes from a cocycle of the procyclic unramified
quotient. Such a cocycle is determined by its value at
Frobenius, modulo the coboundaries
$(q_w^{-e}-1)A_e$. Multiplication by the nonzero
element $q_w^{-e}-1$ of $\mathbb Z_p$ is surjective on $A_e$.
After subtracting a coboundary the cocycle vanishes at
Frobenius, hence on its dense cyclic subgroup, hence
everywhere by continuity. Thus the unramified local
cohomology group is zero.

The same argument applies to $A_{e,\theta}$ at
$\ell\nmid pN$. Its Frobenius scalar is
$\theta(\ell)^{-1}\ell^{-e}$, which is not one: otherwise
the rational number $\ell^e>1$ would be a root of unity.
Again multiplication by the scalar minus one is
surjective on $K/\mathcal O$. This proves both assertions.
At real places $H^1$ is zero because $p$ is odd and the
decomposition group has order two.
\end{proof}

In particular (FC2) has the additional zero local conditions
used for the character groups in the submitted construction.
The use of $e\ge1$ in Lemma~\ref{fcmp-local} is essential.
At $e=0$ the Frobenius scalar can equal one.

\subsection{Shapiro's lemma and integral Fourier maps}

Put $G=G_{\mathbb Q}$, $H=G_L$, and write $\bar g$ for
the image of $g$ in $\Delta$. Consider
\[
 \mathcal M=\operatorname{Map}(\Delta,K/\mathcal O),\qquad
 (g f)(\delta)=\varepsilon_p(g)^{-e}
                                      f(\bar g^{-1}\delta).
                                                               \tag{FC3}
\]
Shapiro's lemma gives
\[
 H^1(G,\mathcal M)\simeq
 H^1(H,(K/\mathcal O)(-e)).
                                                               \tag{FC4}
\]
For clarity, the full cocycle calculation is as follows.
The coinduced module consists of functions $F:G\to K/\mathcal O$
with $F(hx)=\varepsilon_p(h)^{-e}F(x)$ and action
$(gF)(x)=F(xg)$. The identification with (FC3) is
$F(x)=\varepsilon_p(x)^{-e}f(\bar x^{-1})$.
Choose representatives $s(x)$ for the finite right coset
space $H\backslash G$, with $s(H)=1$, and set
\[
 h(x,g)=s(x)g\,s(xg)^{-1}\in H.
                                                               \tag{FC5}
\]
Given a cocycle $c:H\to K/\mathcal O$, define
$C(g)(s(x))=c(h(x,g))$ and extend by the coinduction rule.
The identity
$h(x,g_1g_2)=h(x,g_1)h(xg_1,g_2)$ verifies the cocycle
identity. In the reverse direction use $c(h)=C(h)(1)$.
For an arbitrary $C$, the difference from the reconstruction
of this $c$ is the coboundary of the function
$U(s(x))=C(s(x))(1)$: the cocycle identity gives
\[
 C(g)(s(x))=c(h(x,g))
       +\varepsilon_p(h(x,g))^{-e}U(s(xg))-U(s(x)).
                                                               \tag{FC6}
\]
The formulas preserve coboundaries and continuity, proving
(FC4).

For a rational place $\ell$, the restriction of the coinduced
module to a decomposition group is the direct sum of the
coinduced modules from the decomposition groups at places
of $L$ above $\ell$. This follows by separating the finitely
many orbits on $G/H$; the formulas (FC5)--(FC6) on each
orbit give the corresponding local assertion. Consequently
(FC4) identifies classes which are zero at every finite
decomposition group on its two sides. By
Lemma~\ref{fcmp-local}, this is the coefficient extension of
$\mathcal S_p(L,e)$ on the right.

There are integral, equivariant maps
\[
 \begin{aligned}
 \mathcal F:\mathcal M&\longrightarrow
                         \bigoplus_{\theta\in\widehat\Delta}A_{e,\theta},
 &(\mathcal F f)_\theta
       &=\sum_{\delta\in\Delta}\theta(\delta)^{-1}f(\delta),\\
 \mathcal G:\bigoplus_{\theta\in\widehat\Delta}A_{e,\theta}
                 &\longrightarrow\mathcal M,
 &(\mathcal G a)(\delta)&=\sum_\theta\theta(\delta)a_\theta.
 \end{aligned}                                                 \tag{FC7}
\]
All coefficients are units of $\mathcal O$. The change of
variable $\delta=\bar g\gamma$ shows that the $\theta$
coordinate transforms by
$\theta(\bar g)^{-1}\varepsilon_p(g)^{-e}$, exactly the
character in (FC2). Character orthogonality gives
\[
                    \mathcal G\mathcal F=d\,\mathrm{id},
             \qquad \mathcal F\mathcal G=d\,\mathrm{id}.
                                                               \tag{FC8}
\]
For example, the sum of a nontrivial character is zero
because translation multiplies the sum by a value different
from one; over $K$ this forces the sum to be zero. The
identities in (FC8) are therefore identities over
$\mathcal O$ and on its quotient $K/\mathcal O$.

The action of $j$ under (FC4) corresponds to
$T_jf(\delta)=f(\delta j)$. To check the coefficient sign,
choose a complex conjugation $c\in G$ above $j$.
For a cocycle $C$ with values in (FC3), put
$c_H(h)=C(h)(1)$. The cocycle identity implies
\[
 C(h)(j)=\varepsilon_p(c)^{-e}c_H(chc^{-1})
                      +(1-\varepsilon_p(h)^{-e})C(c)(j).
                                                               \tag{FC9}
\]
The last term is a coboundary. The first is the usual
action of complex conjugation on $H^1(H,(K/\mathcal O)(-e))$.
On the $\theta$ coordinate of (FC7), $T_j$ acts by
$\theta(j)$. Thus the required characters are precisely
\[
                 X_e(L)=\{\theta\in\widehat\Delta:
                                      \theta(j)=(-1)^e\}.
                                                               \tag{FC10}
\]
The factor $(-1)^e$ has not been lost in changing the
coefficient module or the coinduction convention.

\begin{proposition}\label{fcmp-fourier}
Let $\nu=v_p(d)$.
If $\mathcal V_p(e,\theta)$ is finite for every
$\theta\in X_e(L)$, then $\mathcal S_p(L,e)^{j=(-1)^e}$
is finite and
\[
 a_p(L,e)\le\sum_{\theta\in X_e(L)}\ell_p(e,\theta)
             +\operatorname{length}_{\mathbb Z_p}
                             \mathcal S_p(L,e)[p^\nu].
                                                               \tag{FC11}
\]
When $p\nmid d$, the last term is zero and equality holds
between $a_p(L,e)$ and the character sum.
\end{proposition}
\begin{proof}
Since $p$ is odd, the two projectors $(1\pm T_j)/2$
are integral. Apply (FC7) to the image of the projector
$(1+(-1)^eT_j)/2$, and then to locally trivial cohomology.
The resulting map lands in
$\bigoplus_{\theta\in X_e(L)}\mathcal V_p(e,\theta)$.
By (FC8), its kernel is killed by $d$, equivalently by
$p^\nu$. Under (FC4), it is therefore contained in the
$p^\nu$-torsion of $\mathcal S_p(L,e)\otimes_{\mathbb Z_p}\mathcal O$.
When $\nu>0$, finiteness of this torsion, independently
of finiteness of the whole Selmer group, is proved in
Lemma~\ref{fcmp-torsion} below. Kernel and image are finite,
so the stated eigenspace is finite.

Coefficient extension commutes with cochains, cohomology
and the local kernels because $\mathcal O$ is finite free
over $\mathbb Z_p$. It also commutes with the two projectors
and with $p^\nu$-torsion. For a finite $\mathbb Z_p$-module
$T$,
\[
 \operatorname{length}_{\mathcal O}
                  (T\otimes_{\mathbb Z_p}\mathcal O)
                =b\,\operatorname{length}_{\mathbb Z_p}T.
                                                               \tag{FC12}
\]
This follows first from $\mathcal O/p^a\mathcal O$ and
then from the elementary-divisor decomposition of $T$.
Taking lengths and dividing by $b$ proves (FC11).
If $p\nmid d$, (FC8) makes the two Fourier maps inverse
after division by the unit $d$. They preserve every
local condition and the parity projector, which proves
the equality.
\end{proof}

\begin{remark}\label{fcmp-integral-error}
An exact integral decomposition cannot be substituted for
(FC8) when $p\mid d$. For instance, take $\Delta=\mathbb Z/p\mathbb Z$,
let $K$ contain a primitive root $\zeta_p$, and put
$\pi=\zeta_p-1$. In $\operatorname{Map}(\Delta,K/\mathcal O)$
let $f(0)=1/\pi$, $f(1)=-1/\pi$, and let all other
coordinates be zero. This is a nonzero vector, but its
Fourier coefficient at the character $a\mapsto\zeta_p^{ka}$
is $(1-\zeta_p^{-k})/\pi\in\mathcal O$, hence is zero in
$K/\mathcal O$. The same example occurs in either parity
summand for $\mathbb Z/2\mathbb Z\times\mathbb Z/p\mathbb Z$
by applying the integral projector for the first factor.
This explains the necessity of controlling the kernel;
it is not a counterexample to (FC11).
\end{remark}

\subsection{A torsion bound independent of the twist}

\begin{lemma}\label{fcmp-torsion}
Let $L$ be any number field, let $p$ be odd, let $e\ge1$,
and let $\nu\ge1$. Set
$F=L(\mu_{p^\nu})$, $u=[F:L]$, $n_F=[F:\mathbb Q]$,
and let $s_F$ be the number of places of $F$ above $p$.
Then
\[
 \operatorname{length}_{\mathbb Z_p}\mathcal S_p(L,e)[p^\nu]
       \le \nu u+\nu(n_F+s_F)+\frac{\log h_F}{\log p}
       \le \nu u+2\nu n_F+\frac{\log h_F}{\log p}.
                                                               \tag{FC13}
\]
In particular the torsion on the left is finite. The
bound is independent of $e$.
\end{lemma}
\begin{proof}
Put $B=A_e[p^\nu]$. The coefficient exact sequence
$0\to B\to A_e\xrightarrow{p^\nu}A_e\to0$ shows that
every class in $\mathcal S_p(L,e)[p^\nu]$ lifts to
$H^1(L,B)$. At a place away from $p$, inertia acts
trivially on $A_e$ and multiplication by $p^\nu$ on
$H^0(I_w,A_e)=A_e$ is surjective. Hence
$H^1(I_w,B)\to H^1(I_w,A_e)$ is injective.
Every such lift is consequently unramified away from $p$.
It suffices to bound this larger finite-coefficient
cohomology group; no choice of compatible local
primitives above $p$ is required.

Over $F$, the group $B$ is a trivial cyclic group of
order $p^\nu$. Inflation--restriction has kernel
$H^1(\operatorname{Gal}(F/L),B)$, whose cardinality is at
most $p^{\nu u}$: its cocycles are a subgroup of the
functions from a group of order $u$ to $B$. The image
consists of homomorphisms $G_F\to\mathbb Z/p^\nu\mathbb Z$
unramified away from $p$.

The ray class exact sequence and global reciprocity give
\[
 \#\operatorname{Hom}_{\mathrm{unr}\ {\rm outside}\ p}
                   (G_F,\mathbb Z/p^\nu\mathbb Z)
 \le h_F\prod_{w\mid p}
       \#\operatorname{Hom}_{\mathrm{cont}}
                 (\mathcal O_{F_w}^{\times},\mathbb Z/p^\nu\mathbb Z).
                                                               \tag{FC14}
\]
Here is the precise passage. Every homomorphism on the
left factors through a ray class group with modulus
supported above $p$, by reciprocity. The kernel of the
map from that ray class group to the ideal class group
is a quotient of the product of the corresponding
local unit quotients. The number of homomorphisms from
the class group is at most $h_F$; restriction to that
kernel bounds the remaining choices by the displayed
local factors. The groups of homomorphisms for increasing
moduli form an increasing union. The uniform finite
bound just obtained also bounds this union.
These are the ray class exact sequence and reciprocity
in Milne, Chapter~V, Theorems~1.7 and~3.5--3.6.
\cite{milne-cft}

For a finite extension $K_0/\mathbb Q_p$ of degree $n_0$,
\[
 \#\operatorname{Hom}_{\mathrm{cont}}
        (\mathcal O_{K_0}^{\times},\mathbb Z/p^\nu\mathbb Z)
                         \le p^{\nu(n_0+1)}.                   \tag{FC15}
\]
To justify this bound, the residue-field unit group has
order prime to $p$ and its Teichm\"uller representatives
split it off. The principal units form an abelian pro-$p$
group. On a sufficiently deep principal-unit subgroup,
the convergent logarithm and exponential give inverse
isomorphisms with an ideal of $\mathcal O_{K_0}$: choose
the depth to have $p$-adic valuation greater than
$1/(p-1)$, and the usual series converge and their
formal inverse identities hold. This subgroup is
$\mathbb Z_p^{n_0}$ and has finite index. Thus the
principal units are a finitely generated
$\mathbb Z_p$-module of rank $n_0$. Their torsion is a
finite subgroup of a field's multiplicative group,
and is therefore cyclic. The structure theorem over
$\mathbb Z_p$ gives $\mathbb Z_p^{n_0}$ plus one finite
cyclic $p$-group. Each factor contributes at most
$p^\nu$ characters, proving (FC15).

Now $\sum_{w\mid p}[F_w:\mathbb Q_p]=n_F$.
Equations (FC14)--(FC15) give the bound
$h_Fp^{\nu(n_F+s_F)}$. Combining with the
inflation--restriction kernel and taking $\log_p$
proves the first inequality in (FC13).
The second uses $s_F\le n_F$.
\end{proof}

\subsection{The growing family and its uniform error}

Return to the fields
$L=E_h=K_\chi R_h(\sqrt{-q},\sqrt{-d_0})$
of (IV1), writing $d_0$ for the second auxiliary prime
to distinguish it from the degree $d=[L:\mathbb Q]$.
The fixed primitive character $\chi$ has conductor $M$
and order $o$. The degree of $K_\chi$ is $o$ and that of
$R_h$ is $2^h$; the two quadratic adjunctions show
\[
                            d\mid 2^{h+2}o.
                                                               \tag{FC16}
\]
For each odd prime $p\mid o$, put
\[
 a_p^0=v_p(o),\qquad m_p=\varphi(p^{a_p^0}),\qquad
 C_p=3a_p^0m_p\frac{\log p}{\log2}
       +\frac{m_p^2}{2}
                    \left(1+a_p^0\frac{\log p}{\log2}\right),
 \quad C_o=\sum_{\substack{p\mid o\\p\ne2}}C_p.
                                                               \tag{FC17}
\]
These constants depend only on the original character.
In particular $C_o=0$ if $o$ is a power of two.

\begin{theorem}\label{fcmp-uniform}
Let $P$ be a finite set of odd primes and $t\ge2$ an integer.
Suppose that the groups in (FC2) are finite for every
$p\in P$, $1\le e<t$, and $\theta\in X_e(E_h)$. Then
\[
 \begin{split}
 \sum_{p\in P}\log p\sum_{e=1}^{t-1}a_p(E_h,e)
 &\le
 \sum_{p\in P}\log p\sum_{e=1}^{t-1}
                    \sum_{\theta\in X_e(E_h)}\ell_p(e,\theta)\\
 &\quad+C_o N_h^2(t-1)\log N_h .
 \end{split}                                                   \tag{FC18}
\]
The coefficient of each character length is exactly one,
regardless of the coefficient extension. The error has
a fixed polynomial bound in $N_ht$.
\end{theorem}
\begin{proof}
By (FC16), the Fourier error is zero unless $p\mid o$,
and otherwise $\nu=v_p(d)\le a_p^0$. In (FC13) one may
enlarge $F$ to $E_h(\mu_{p^{a_p^0}})$; the same proof
uses $p^\nu$ coefficients and gives the same formula
with this larger field. Then
\[
 u\le m_p,\quad n_F\le dm_p,\quad
 F\subset\mathbb Q(\mu_{N_hp^{a_p^0}}),\quad
 |\operatorname{disc}F|
                 \le (N_hp^{a_p^0})^{n_F}.                     \tag{FC19}
\]
The discriminant estimate follows from the tower formula
and the cyclotomic derivative-norm argument in the proof
of Lemma~\ref{sv-field}.
Corollary~\ref{sl-class-number} and $\log(1+x)\le x$ give
$\log h_F\le \tfrac12 n_F\log|\operatorname{disc}F|$.
Consequently the error in (FC11), multiplied by $\log p$,
is at most
\[
 a_p^0m_p\log p+2a_p^0dm_p\log p
            +\frac12d^2m_p^2(\log N_h+a_p^0\log p)
 \le C_p d^2\log N_h .
                                                               \tag{FC20}
\]
For the last inequality use $d\ge1$ and $\log N_h\ge\log2$.
Sum (FC20) over the fixed primes dividing $o$ and over
$1\le e<t$, and use $d\le N_h$. Together with (FC11)
this proves (FC18). Since $\log N_h\le N_h$ and $C_o$
is fixed, its error is bounded by $(N_ht)^C$ for one
integer $C$ chosen before $h,q,d_0,t$.
\end{proof}

In the source notation the characters in (FC10) are
exactly the parity-restricted characters of the chosen
abelian field, and (FC2) is the normalized length of
the group with zero local conditions at primes dividing
$pN_h$.\cite{alpo-uniform-source}
Thus (FC18) supplies the whole-field comparison needed
to assemble the complementary-parity estimates.
It does not require that $p$ be prime to the field
degree. In particular a bound $CB$ for the total
fixed-conductor character contribution remains $CB$:
the passage to $E_h$ adds only the polynomial error
in (FC18). The unproved character estimates in the
remaining automorphic cases are still required.


